% Part: lambda-calculus
% Chapter: syntax
% Section: abbreviated-syntax

\documentclass[../../../include/open-logic-section]{subfiles}

\begin{document}

\olfileid{lam}{syn}{abb}
\olsection{संक्षिप्त विन्यास}

\olref[trm]{defn:term} मध्ये परिभाषित पदे लिहिताना
कधीकधी अवजड वाटतात; म्हणून अधिक संक्षिप्त विन्यास
मांडणे उपयुक्त आहे. अर्थात, या संक्षिप्त संकेतनातील
पदेही एकमेव रीतीने वाचता येतील याची काळजी
घ्यावी लागेल. \emph{संक्षिप्त पदे} म्हणून ओळखली
जाणारी एक प्रचलित पद्धत पुढीलप्रमाणे आहे.

\begin{enumerate}
\item कंस वगळले असता उपयोजन डावीकडून उजवीकडे
  होते. उदाहरणार्थ, $M$, $N$, $P$ आणि~$Q$
  ही पदे असतील, तर $MNPQ$ हे $(((MN)P)Q)$
  चे संक्षिप्त रूप आहे.
\item पुन्हा, कंस वगळले असता लॅम्डा अमूर्तीकरणाला
  शक्य तितकी व्यापक व्याप्ती द्यायची. उदाहरणार्थ,
  $\lambd[x][MNP]$ याचा अर्थ
  $(\lambd[x][MNP])$ असा घ्यायचा.
\item एका लॅम्डाने अनेक चरे अमूर्त करता येतात.
  उदाहरणार्थ, $\lambd[xyz][M]$ हे
  $\lambd[x][\lambd[y][\lambd[z][M]]]$
  चे संक्षिप्त रूप आहे.
\end{enumerate}

उदाहरणार्थ,
\[
\lambd[xy][xxyx \lambd[z][xz]]
\]
हे पुढीलचे संक्षिप्त रूप आहे:
\[
(\lambd[x][(\lambd[y][((((xx)y)x)(\lambd[z][(xz)]))])]).
\]

\begin{prob}
संक्षिप्त पद
$\lambd[g][(\lambd[x][g (x x)])
  \lambd[x][g (x x)]]$
विस्तारून लिहा.
\end{prob}

\end{document}
